% generated by scripts/paper_framing.py, rule framing-v2 (chosen after reading, see DECISIONS 2026-11-23T20:30Z); do not edit
\paragraph{Main findings.} On an adapted $500$-graph VGCure evaluation (drawings by other authors in a style none of our adapters saw, with our prompts, images scaled to our input size and strict scoring), the repair-trained adapter (arm~D) answers edge-existence and reachability questions better than direct question-answer fine-tuning (arm~B): the difference is $+0.056$ (95\% interval [$+0.038$, $+0.075$], mean of 3 seeds, above zero at every seed). The lead comes from both tasks: edge existence $+0.049$ and reachability $+0.063$. On our own pre-registered locked test, which decides the hypothesis, D does worse than B: D$-$B is $-0.105$ (95\% interval [$-0.117$, $-0.093$]), so the hypothesis is not supported at this scale. Section~\ref{sec:why} traces the disagreement: on our test D's deficit is in the critical-edge task, where it answers \emph{true} to 79\% of the questions, while its reachability is level with B's; on VGCure, which has no critical-edge items, D gains on false edge claims and unreachable pairs. This is a pilot-scale study: all adapters were trained on $512$ synthetic root graphs with one 2B base model (Qwen3.5-2B, LoRA), and nothing here measures real diagrams. Both evaluations were fixed before any of their outputs existed; giving the external result this place in the abstract and introduction was decided after both were read, while the wording rule fixed in advance put only the locked-test statistic here (Section~\ref{sec:why}). Per-split, per-task and secondary comparisons are in Section~\ref{sec:test-results}.
